\section{伪全纯截面}\label{sec:17}
\subsection{截面计数}\label{sec:17a}

推广 Floer TQFT 标准框架的一种方式，是将目标空间换成 Lefschetz 纤维化，并以伪全纯截面代替映射。先在最简单的情形，即底曲面紧致时解释这一推广，再扩展框架以允许条带状端。这样做的目的，是将代数扭转构造 $T$ 与几何 Dehn 扭转 $\tau$ 联系起来，见\stref{推论}{17.17}。我们会特别仔细地以相当典范的形式建立这种认同，使之易于推广到 $\Z/2$ 等变背景。

令 $\pi:E^{2n+2}\to S$ 为 Lefschetz 纤维化，仍照例正合。\emph{Lagrangian 边界条件}是一个包含在 $\partial^vE$ 中且与 $\partial^hE$ 不交的子流形 $F^{n+1}\subset E$，并要求存在 $\alpha_F\in\Omega^1(\partial S)$ 及 $h_F\in C^\infty(F,\R)$，使
\begin{equation}\label{eq:17.1}
\theta_E|F=(\pi|F)^*(\alpha_F+\theta_S|\partial S)+dh_F.
\end{equation}
稍加解释。$F$ 显然在 $E$ 中为 Lagrangian。由于 $\pi$ 在 $\partial^vE$ 中没有临界点，$\pi|F:F\to\partial S$ 必为纤维丛，见\stref{引理}{16.2}。每个 $F_z\subset E_z$，$z\in\partial S$，实际上都是正合 Lagrangian 子流形，且这些子流形由沿边界的平行移动互相映到彼此，见\stref{引理}{16.3}。逆命题也在很大程度上成立：假设 $F_z\subset E_z$，$z\in\partial S$，是一族由边界平行移动互相映到彼此的正合 Lagrangian 子流形，令 $F$ 为它们的并。则 $\theta_E|F$ 闭，且可找到函数 $h_F$，使 $\theta_E|F-dh_F$ 限制到每个 $F_z$ 时消失。若
\SourcePage{244}
再假设 $F_z$ 连通，则 $\theta_E|F-dh_F$ 必是 $\partial S$ 上某个一形式的拉回，故 $F$ 为 Lagrangian 边界条件。\znote{236}{“正合 Lagrangian 边界条件”的准确含义是这个分解，而不只是“$F$ 为 Lagrangian”。它由三件事合作给出：$\pi$ 在 $\partial^vE$ 上无临界点，故 $\pi|F$ 是纤维丛（引理 16.2）；每条纤维 $F_z$ 正合、且由平行移动彼此相联（引理 16.3）；再加上 $F_z$ 连通，使拼出的 $\theta_E|F-dh_F$ 只能来自 $\partial S$ 上的一形式。少了连通性，各纤维上的 $h_F$ 之间可差局部常数，$F$ 就未必满足\eref{17.1}。[来源：原书 §(17a)，印刷页 236。]}

固定这样的 $F$。\emph{相对扰动数据}是一对 $(K,J)$，其中：$K$ 为 $E$ 上一形式，在 $TE^v$ 中的向量上为零，在 $\partial^hE\cup\operatorname{Crit}(\pi)$ 的某个邻域中完全消失，且满足 $K|F=0\in\Omega^1(F)$；$J$ 为 $E$ 上的 $\omega_E$ 相容殆复结构，使 $\pi$ 伪全纯，并在 $\partial^hE\cup\operatorname{Crit}(\pi)$ 的某个邻域中与 $I_E$ 一致。给定这些数据，边界位于 $F$ 中的\emph{非齐次伪全纯截面}是映射 $u:S\to E$，满足
\begin{equation}\label{eq:17.2}
\begin{cases}
\pi(u(z))=z,\\
Du(z)+J(u)\circ Du(z)\circ I_S=Y(u)+J(u)\circ Y(u)\circ I_S,\\
u(\partial S)\subset F.
\end{cases}
\end{equation}
这里 $Y$ 是丛 $\Hom(\pi^*TS,TE^v)$ 的截面，按下述方式由 $K$ 确定：给定任意 $\xi\in TS_z$，$z=\pi(x)$，在 $E_z$ 的正则部分上将其提升为切向量场 $X$；函数 $K(X)$ 与提升的选择无关，定义 $Y(\xi)_x$ 为该函数在 $E_z$ 上关联的 Hamilton 向量场于 $x$ 处的值。由于要求 $K$ 在临界点附近为零，$Y$ 也如此，故可平凡地延伸过这些点。

$Y$ 在临界点附近消失，因而可连续延伸过这些点；在 $\partial^hE$ 附近也消失。

\StatementNumber{1}\begin{note}\label{st:17.1}
相对扰动数据有一个简单的几何解释。考虑修改后的二形式
\begin{equation}\label{eq:17.3}
\omega_E^{\mathrm{new}}=\omega_E-dK,
\end{equation}
它仍闭且逐纤维为辛；其关联联络为
\begin{equation}\label{eq:17.4}
TE_x^{h,\mathrm{new}}=\{X=X^h+X^v\in TE_x=TE_x^h\oplus TE_x^v:
X^v=Y(D\pi(X))\}.
\end{equation}
关于殆复结构，注意按假设 $J$ 保持分解 $TE^h\oplus TE^v$，且在第一个直和项上等于 $I_E$。故这里的选择实质上是 $TE^v$ 上殆复结构的选择。存在唯一的 $J^{\mathrm{new}}$，使 $\omega_E^{\mathrm{new}}(\cdot,J^{\mathrm{new}}\cdot)$ 对称、$\pi$ 为 $J^{\mathrm{new}}$ 全纯，且在 $TE^v$ 上有 $J^{\mathrm{new}}=J$。用这些术语，\eref{17.2}就是截面在通常意义下为 $J^{\mathrm{new}}$ 全纯的方程。联络观点也有助于理解解 $u$ 的一些基本性质。例如，能量的先验估计为
\[
E(u)=\int_S\tfrac12|(Du)^v-Y(u)|^2
=\int_{\partial S}(\alpha_F+\theta_S)-\int_Su^*R,
\]
\SourcePage{245}
其中 $(\cdot)^v$ 是关于 $TE=TE^h\oplus TE^v$ 的垂直分量，而曲率项在局部平凡化\eref{15.3}中由下式给出：
\begin{equation}\label{eq:17.5}
\begin{aligned}
R(\partial_s,\partial_t)=\bigl(&\partial_s(K-\kappa)(\partial_t)-\partial_t(K-\kappa)(\partial_s)\\
&+\{(K-\kappa)(\partial_s),(K-\kappa)(\partial_t)\}\bigr)\,ds\wedge dt.
\end{aligned}
\end{equation}
等价地，使用修改后的联络\eref{17.4}，可写为 $E(u)=\int_S\tfrac12|(Du)^{v,\mathrm{new}}|^2$；此时\eref{17.5}确实是该联络的曲率。\znote{237}{$R$ 的表达式值得逐项读：前两项是通常的 $d(K-\kappa)$，第三项 $\{(K-\kappa)(\partial_s),(K-\kappa)(\partial_t)\}$ 来自 $K-\kappa$ 取值为垂直向量场、取值之间不可交换的二次修正；漏掉它，$R$ 就不闭、能量恒等式也不成立。能量恒等式的意义是：非齐次项 $Y$ 只贡献垂直方向的能量，而右端把它换成边界项减去曲率积分。于是 $R$ 有界时能量受边界项控制，这正是下文紧性论证的出发点；也解释了为何要求 $K$ 在 $\partial^hE\cup\operatorname{Crit}(\pi)$ 附近为零——否则 $R$ 在那里没有一致界。[来源：原书 §(17a)，印刷页 237。]}
\end{note}

\StatementNumber{2}\begin{note}\label{st:17.2}
在此背景下，也有移动边界条件的对应概念。这种边界条件只是任意一族正合 Lagrangian 子流形 $F_z\subset E_z$，$z\in\partial S$。适当的相对扰动数据概念是一对 $(K,J)$，其中不再要求 $K$ 在 $F$ 上的限制为零，而要求：对任意 $X\in TF|F_z$，$z\in\partial S$，若 $D\pi(X)$ 平行于边界，则
\begin{equation}\label{eq:17.6}
d(i_XK)|F_z=\omega_E(\cdot,X)|F_z.
\end{equation}
除此之外，非齐次伪全纯截面的定义不变。其含义应该很清楚：$F_z$ 未必在 $E$ 上给定辛联络的平行移动下保持，但对于满足\eref{17.6}的某个 $K$，转到修改后的形式\eref{17.3}，就恰好恢复了这种相容性。
\end{note}

记 $\M_{E/S}$ 为\eref{17.2}所有解的模空间。在其中任意点 $u$ 处，该方程的线性化由向量丛 $u^*(TE^v)$ 上的 $\bar\partial$ 算子 $D_{E/S,u}$ 给出，边界值在 $\partial S$ 上由 $u^*(TF^v)$ 给定。如通常一样，若 $D_{E/S,u}$ 对所有 $u$ 都满射，则称 $(K,J)$ 正则；此时模空间为闭光滑流形。一般选择 $(K,J)$ 就具有这一性质。取任意这样的选择，可通过对孤立点作模二计数，定义一个简单模型不变量：
\begin{equation}\label{eq:17.7}
\Phi_{E/S}=\#\M_{E/S}\in\Z/2.
\end{equation}

\clearpage
\StatementNumber{3}\begin{example}\label{st:17.3}
令 $\pi:E\to S$ 为\stref{例题}{15.4}和\stref{例题}{16.5}中的局部模型。赋予它 Lagrangian 边界条件 $F$，取为 $|z|=r$ 时标准消失圈 $F_z=\sqrt zS^n$ 之并。在 $F$ 上有 $x_j=(\pi(x)/r)\bar x_j$，故
\[
\theta_E|F=\frac i4\sum_j(x_j\,d\bar x_j-\bar x_j\,dx_j)
=\pi^*\left(\frac{i}{8r}(z\,d\bar z-\bar z\,dz)\right).
\]
选取 $\alpha_F$，使 $\alpha_F+\theta_S|\partial S$ 对应于右边的形式，并取 $h_F=0$，就表明 $F$ 是 Lagrangian 边界条件。当然，这一点此前已经知道，因为 $\sqrt zS^n$ 在任意方向的平行移动下都保持。
\SourcePage{246}
取标准复结构 $J=I_E$ 和平凡非齐次项 $K=0$。所得模空间 $\M_{E/S}$ 可显式确定：它由映射
\begin{equation}\label{eq:17.8}
u_a:S\longrightarrow E\subset\CC^{n+1},\qquad
u_a(z)=r^{-1/2}az+r^{1/2}\bar a
\end{equation}
组成，其中 $a\in\CC^{n+1}$ 满足 $a_1^2+\cdots+a_{n+1}^2=0$ 及 $\|a\|^2=1/2$，见\bcite{131}引理 2.16。此外，这些全都是模空间的正则点。严格地说，还须稍加小心，即\eref{15.10}中的 $s$ 应充分大，使所有这些截面的像确实包含在 $E$ 中。这与 $\partial^hE$ 并非全纯凸有关，也正是这一点使 $E$ 不能成为完全有效的正合 Lefschetz 纤维化例子。暂且忽略这一缺陷，就会发现不变量\eref{17.7}因维数理由而平凡，因为模空间没有孤立点。
\end{example}

这一计算是一个更一般消失结果的模型。令 $\pi:E\to S$ 为底空间 $S$ 是圆盘、且只有一个临界点的 Lefschetz 纤维化。选择 $\partial S$ 上一点 $z$，及从 $z$ 到临界值的消失路径 $\gamma$。记 $M=E_z$、$V=V_\gamma$。由\eref{16.6}，沿边界回路 $\lambda$ 的单值变换 $h_\lambda$ 与 $\tau_V$ 同痕，且同痕在保持端点不动的同伦意义下唯一。这意味着，可以如下为纤维化赋予移动 Lagrangian 边界条件 $F$：从 $F_z=V$ 出发，通过平行移动绕边界几乎一周，再用同痕 $h_\lambda(V)\simeq\tau_V(V)=V$ 将结果接回 $F_z$。称之为\emph{标准边界条件}。

\StatementNumber{4}\begin{lemma}\label{st:17.4}
$\Phi_{E/S}=0$。
\end{lemma}
\begin{proof}
这里只概述证明；细节见\bcite{131}第 2.3 节。为简化记号，假设 $S=D$ 为单位圆盘，且 $\operatorname{Critv}(\pi)=\{0\}$。作不影响不变量的局部变形后，可假设临界点附近存在全纯 Morse 坐标，辛形式在其中标准化。考虑将给定纤维化限制到更小圆盘 $S^r=rD$ 所得的 Lefschetz 纤维化族 $\pi^r:E^r\to S^r$，$r\in(0,1]$。从 $F$ 出发，可为它们赋予光滑边界条件族 $F^r$，在 $r$ 小时等于\stref{例题}{17.3}局部模型中的边界条件。为 $(E^r,F^r)$ 取一族扰动数据，在 $r$ 小时平凡。此时，紧致性论证\bcite{131}引理 2.15 表明，在极限 $r\to0$ 中，$\M_{E^r/S^r}$ 的所有元素都包含在临界点任意小的指定邻域内，因而回到\stref{例题}{17.3}的情形。最后，用常规配边论证将结果带回原来的 $(E,F)$。
\end{proof}

\SourcePage{247}
作为补充，再提及一个较初等但普遍性较弱的证明，它只依赖于模空间的维数。先需要关于 $(E,F)$ 拓扑的一个基本事实：

\StatementNumber{5}\begin{lemma}\label{st:17.5}
若相对第一 Chern 类 $2c_1(M,V)\in H^2(M,V)$ 消失，则相应的 $2c_1(E,F)\in H^2(E,F)$ 等于 $H^2(S,\partial S)$ 标准生成元的拉回的 $(n+1)$ 倍。
\end{lemma}
\begin{proof}
由\secref{15c}可知，$E$ 可赋予二次复体积形式，故 $2c_1(E)=0$，这蕴含 $2c_1(E,F)$ 来自 $H^1(F)$ 的元素。对 $n>1$，这意味着它是 $H^2(S,\partial S)$ 某个元素在 $\pi^*$ 下的像。对 $n=1$，注意 $2c_1(E,F)$ 到 $H^2(\Delta,V)$ 的限制必因 $\Delta$ 的 Lagrangian 性质而消失，也能得出同样结论；这里 $\Delta$ 为 Lefschetz 顶针。有了这一认识，论证便归结为局部模型中的计算，这里省略。
\end{proof}

模空间 $\M_{E/S}$ 在点 $u$ 处的预期维数由算子 $D_{E/S,u}$ 的指标给出。应用相对 Riemann–Roch 公式，即稍作改写的\stref{引理}{11.7}，及\stref{引理}{17.5}，得到
\[
\operatorname{index}(D_{E/S,u})
=n+\int_S2c_1(u^*(TE^v),u^*(TF^v))=2n-1,
\]
这当然与局部模型中看到的一致。特别地，无论选择何种正则相对扰动数据，$\M_{E/S}$ 都不会有孤立点。这与前一个更一般的配边论证所得结果有所不同；前者只保证这些孤立点的贡献相互抵消。\znote{239}{指标公式给出的是模空间的“形式维数” $2n-1$，与例题 17.3 局部模型中看到的维数一致；要点是 $2n-1>0$，故正则情形下根本没有孤立点，$\Phi_{E/S}=0$ 便因维数理由自动成立。两种论证的差别正在此处：配边论证只说明孤立点的贡献成对抵消，维数论证则说明孤立点压根不存在。真正的困难始终是 $\M_{E/S}$ 的紧性——紧性失效时，形式维数与实际维数可以不符，孤立点仍可能出现，这才是引理 17.4 必须单独证明的原因。[来源：原书 §(17a)，印刷页 239。]}
